% page 288 of the facsimile (image p-287 of the scan)
% block 152.4 x 237.7 mm ; left margin 8.0 mm
\pgc{-12.700mm}{0.000mm}{
\l{\cel{3}280}
\l{}
\l{\sou{klinometro}. Aparato uzata - en aernavigado,marnavigado-}
\l{\cel{3}(an busolo) por mezurar la deviaco de la horizontalo.}
\l{\cel{3}- DFRS.}
\l{}
\l{\sou{klipero}. (\sou{navig.}) Seglo-navo, multa-tuna, pasable rapida,}
\l{\cel{3}uzata olim por la teo-komerco inter Britania e Chinia.}
\l{\cel{3}- DEFIRS.}
\l{}
\l{\sou{klishar}. (\sou{trans.}) I. (\sou{imprim.}) Reproduktar reliefe la traco}
\l{\cel{3}de kompostajo ek tipi movebla, de la tipo originala di}
\l{\cel{3}graburo, e c.) - II. (\sou{fctogr.}) Facar la vitro-plako ne-}
\l{\cel{3}gativa, de qua la fotografisto obtenos la pozitivo.-}
\l{\cel{3}III. (\sou{hektogr}.) (a) Perforar la folio de materio vaxatra}
\l{\cel{3}(per la tipi di skribo-mashino o per stilo se koncern-}
\l{\cel{3}esas manuo-skriburo o desegnuro) qua uzesos por la hekto-}
\l{\cel{3}grafilo ; (b) skribar (o desegnar) sur la folio de ma-}
\l{\cel{3}terio paperatra, tre glata, per inko specala (hektograf-}
\l{\cel{3}inko), to quo uzesos por la hektografilo. - DFRS.}
\l{}
\l{\sou{klistero}. La medikamento liquida quan onu injektas, tra la}
\l{\cel{3}anuso, aden la grosa intestino. - DEFIRS.}
\l{}
\l{\sou{klitorido}. (\sou{anat.)}\cel{2}Mikra organo karnaja, enire di la}
\l{\cel{3}vulvo. - DEFIS.}
\l{}
\l{\sou{klivar}. (\sou{trans.}) Dividar (diamanti, kristali, mikao-bloko)}
\l{\cel{3}ad strati lamela. - DEF.}
\l{}
\l{\sou{kloako}.\cel{2}I. Loko destinita ad esor kolektuyo di la rezi-}
\l{\cel{3}dui. - II. (\sou{anat.}) Posho quan, en la korpo di la uceli,}
\l{\cel{3}formacas la extremajo di la intestino. - DEFIRS.}
\l{}
\l{\sou{kloko.}\cel{2}La 24-esma parto di dio, determinata per horlo-}
\l{\cel{3}jo. - E.}
\l{}
\l{\sou{klonuso}. Fenomeno qua rezultas de la exajereso di la}
\l{\cel{3}reflexi. - DEFIRS.}
\l{}
\l{\sou{kloro}. (\sou{kemio)}. Korpo simpla, gasa, flava verdatre, odor-}
\l{\cel{3}sufokiva, qua saporas kaustike, uzata kom senkolor-}
\l{\cel{3}igivo e kom desinfektivo. (\cel{1}\sou{Simb. kem.}:\cel{2}Chl. - DEFIRS}
\l{}
\l{\sou{klorofilo}. (\sou{bot.}) Materio qua donas a la parti verda di}
\l{\cel{3}la planto olia koloro. - DEFS.}
\l{}
\l{\sou{kloroformo.}\cel{2}Substanco anesteziala, liquida, oleatra,}
\l{\cel{3}senkolora. - DEFIRS.}
\l{}
\l{\sou{kloroleucito}. (\sou{biol.})\cel{2}Granulo de protoplasmo, impregn-}
\l{\cel{3}ita ye klorofilo.}
\l{}
\l{kloroplasto. (biol.) Rolo quan pleas la kloroleucito}
\l{\cel{3}la produkto di klorofilo.}
}
